\documentclass[11pt,a4paper]{article}

\usepackage[indonesian]{babel}
\usepackage{amsmath,amssymb,amsthm}
\usepackage{enumitem}
\usepackage{array}
\usepackage{geometry}
\geometry{
  left=20mm,
  right=20mm,
  top=20mm,
  bottom=25mm
}

\newcommand{\E}{\mathbb{E}}
\newcommand{\PP}{\mathbb{P}}
\newcommand{\R}{\mathbb{R}}
\newcommand{\F}{\mathcal{F}}

\begin{document}
\pagenumbering{gobble}

\begin{center}
  {\Large \textbf{EVALUASI AKHIR SEMESTER GENAP 2025/2026}}\\
  {\Large \textbf{PROGRAM STUDI MAGISTER MATEMATIKA}}\\
  {\Large \textbf{DEPARTEMEN MATEMATIKA FSAD ITS}}
\end{center}

\vspace{0.4cm}

\begin{tabular}{@{}ll@{}}
  Mata kuliah / Kode & : Kalkulus Stokastik, SM235312 \\
  Hari / Tanggal     & : Senin, 15 Juni 2026 \\
  Waktu / Sifat      & : 100 menit / Buka Buku \\
  Kelas / Dosen      & : Dr. mont. Kistosil Fahim, M.Si.
\end{tabular}

\vspace{0.6cm}

\begin{center}
  \textbf{SOAL}
\end{center}

\begin{enumerate}[leftmargin=*, itemsep=0.5cm]
  \item Suatu eksperimen terdiri dari pelemparan dua dadu bersisi delapan. Didefinisikan
        variabel acak $\xi$ dan $\eta$ secara berturut-turut sebagai modulo $5$ dan $4$
        dari jumlah nilai dua dadu yang dilempar. Tentukan ekspektasi bersyarat
        $\E(\xi\mid\eta)$.

  \item Misalkan $\theta\in\R$ dan $(W_t)_{t\ge 0}$ adalah proses Wiener standar pada
        ruang probabilitas $(\Omega,\F,\PP)$ yang dilengkapi dengan filtrasi alami
        $(\F_t)_{t\ge 0}=\sigma(\{W(r):0\le r\le t\})$. Didefinisikan proses
        \[
          M_t=\exp\left(\theta W_t-\frac{\theta^2}{2}t\right),\qquad t\ge 0.
        \]
        \begin{enumerate}[label=\alph*)]
          \item Buktikan bahwa $\E[M_t]=1$ untuk setiap $t\ge 0$.
          \item Tunjukkan bahwa $\E(M_t\mid\F_s)=M_s$ untuk $0\le s<t$.
          \item Apakah $(M_t)_{t\ge 0}$ adalah martingale terhadap filtrasi $(\F_t)_{t\ge 0}$?
                Jelaskan.
        \end{enumerate}

  \item Misal $W(t)$ proses Wiener dan $Y(t)=\ln S(t)$ dengan
        \[
          dS(t)=\mu S(t)\,dt+\sigma S(t)\,dW(t),
        \]
        dengan $\mu,\sigma\in\R$. Dengan menggunakan formula It\^o, buktikan bahwa
        \[
          Y(t)=Y(0)+\left(\mu-\frac{\sigma^2}{2}\right)t+\sigma W_t.
        \]

  \item Misalkan $f$ dan $g$ adalah fungsi dari $\R$ ke $\R$ yang turunan pertamanya
        terbatas, yakni ada bilangan real positif $M$ sehingga untuk setiap $x,y\in\R$
        berlaku
        \[
          |f'(x)|\le M,\qquad |g'(x)|\le M.
        \]
        Lebih lanjut $\xi_0$ adalah variabel acak terintegral kuadrat $\F_0$-terukur.
        Tunjukkan bahwa masalah nilai awal
        \[
          d\xi(t)=f(\xi(t))\,dt+g(\xi(t))\,dW(t),\qquad \xi(0)=\xi_0
        \]
        mempunyai solusi tunggal.
\end{enumerate}

\newpage

\begin{center}
  \textbf{PENYELESAIAN}
\end{center}

\begin{enumerate}[leftmargin=*, itemsep=0.6cm]
  \item Misalkan $X_1$ dan $X_2$ menyatakan hasil dua dadu bersisi delapan, sehingga
        $X_1,X_2\in\{1,2,\ldots,8\}$ dan
        \[
          S:=X_1+X_2\in\{2,3,\ldots,16\}.
        \]
        Karena semua $64$ pasangan hasil lemparan ekuipeluang, banyaknya pasangan yang
        menghasilkan jumlah $s$ adalah
        \[
          c_s=
          \begin{cases}
            s-1, & 2\le s\le 9,\\
            17-s, & 10\le s\le 16.
          \end{cases}
        \]
        Didefinisikan
        \[
          \xi=S \pmod 5,\qquad \eta=S \pmod 4.
        \]

        Untuk setiap nilai $\eta$, kita hitung $\E(\xi\mid \eta=j)$.

        Jika $\eta=0$, maka $S\in\{4,8,12,16\}$, dengan bobot berturut-turut
        $3,7,5,1$, sedangkan nilai $\xi$ berturut-turut $4,3,2,1$. Jadi
        \[
          \E(\xi\mid\eta=0)
          =\frac{3\cdot 4+7\cdot 3+5\cdot 2+1\cdot 1}{3+7+5+1}
          =\frac{44}{16}
          =\frac{11}{4}.
        \]

        Jika $\eta=1$, maka $S\in\{5,9,13\}$, dengan bobot $4,8,4$, dan nilai
        $\xi$ adalah $0,4,3$. Jadi
        \[
          \E(\xi\mid\eta=1)
          =\frac{4\cdot 0+8\cdot 4+4\cdot 3}{4+8+4}
          =\frac{44}{16}
          =\frac{11}{4}.
        \]

        Jika $\eta=2$, maka $S\in\{2,6,10,14\}$, dengan bobot $1,5,7,3$, dan nilai
        $\xi$ adalah $2,1,0,4$. Jadi
        \[
          \E(\xi\mid\eta=2)
          =\frac{1\cdot 2+5\cdot 1+7\cdot 0+3\cdot 4}{1+5+7+3}
          =\frac{19}{16}.
        \]

        Jika $\eta=3$, maka $S\in\{3,7,11,15\}$, dengan bobot $2,6,6,2$, dan nilai
        $\xi$ adalah $3,2,1,0$. Jadi
        \[
          \E(\xi\mid\eta=3)
          =\frac{2\cdot 3+6\cdot 2+6\cdot 1+2\cdot 0}{2+6+6+2}
          =\frac{24}{16}
          =\frac{3}{2}.
        \]

        Dengan demikian,
        \[
          \E(\xi\mid\eta)=
          \begin{cases}
            \dfrac{11}{4}, & \eta=0,\\[1mm]
            \dfrac{11}{4}, & \eta=1,\\[1mm]
            \dfrac{19}{16}, & \eta=2,\\[1mm]
            \dfrac{3}{2}, & \eta=3.
          \end{cases}
        \]
        Atau secara ringkas,
        \[
          \E(\xi\mid\eta)
          =\frac{11}{4}\mathbf{1}_{\{\eta=0\}}
          +\frac{11}{4}\mathbf{1}_{\{\eta=1\}}
          +\frac{19}{16}\mathbf{1}_{\{\eta=2\}}
          +\frac{3}{2}\mathbf{1}_{\{\eta=3\}}.
        \]

  \item \begin{enumerate}[label=\alph*)]
          \item Karena $W_t\sim\mathcal{N}(0,t)$, maka fungsi pembangkit momen peubah normal
                memberi
                \[
                  \E\bigl(e^{\theta W_t}\bigr)=e^{\theta^2 t/2}.
                \]
                Akibatnya
                \[
                  \E[M_t]
                  =\E\left[\exp\left(\theta W_t-\frac{\theta^2}{2}t\right)\right]
                  =e^{-\theta^2 t/2}\E(e^{\theta W_t})
                  =e^{-\theta^2 t/2}e^{\theta^2 t/2}
                  =1.
                \]

          \item Untuk $0\le s<t$, tulis
                \[
                  W_t=W_s+(W_t-W_s).
                \]
                Maka
                \begin{align*}
                  M_t
                   & =\exp\left(\theta W_t-\frac{\theta^2}{2}t\right)\\
                   & =\exp\left(\theta W_s-\frac{\theta^2}{2}s\right)
                  \exp\left(\theta(W_t-W_s)-\frac{\theta^2}{2}(t-s)\right)\\
                   & =M_s\exp\left(\theta(W_t-W_s)-\frac{\theta^2}{2}(t-s)\right).
                \end{align*}
                Karena inkremen Wiener $W_t-W_s$ independen terhadap $\F_s$, maka
                \begin{align*}
                  \E(M_t\mid\F_s)
                   & =M_s\,
                  \E\left(
                  \left.
                  \exp\left(\theta(W_t-W_s)-\frac{\theta^2}{2}(t-s)\right)
                  \right|\F_s
                  \right)\\
                   & =M_s\,
                  \E\left[
                  \exp\left(\theta(W_t-W_s)-\frac{\theta^2}{2}(t-s)\right)
                  \right].
                \end{align*}
                Sekarang $W_t-W_s\sim\mathcal{N}(0,t-s)$, sehingga dengan cara yang sama
                seperti pada butir (a),
                \[
                  \E\left[
                  \exp\left(\theta(W_t-W_s)-\frac{\theta^2}{2}(t-s)\right)
                  \right]
                  =1.
                \]
                Jadi
                \[
                  \E(M_t\mid\F_s)=M_s.
                \]

          \item Ya, $(M_t)_{t\ge 0}$ adalah martingale terhadap filtrasi
                $(\F_t)_{t\ge 0}$, sebab:
                \begin{enumerate}[label=\roman*)]
                  \item $M_t$ adalah $\F_t$-terukur karena merupakan fungsi dari $W_t$,
                  \item $M_t$ terintegralkan karena $\E|M_t|=\E[M_t]=1<\infty$,
                  \item dari butir (b) berlaku $\E(M_t\mid\F_s)=M_s$ untuk setiap $0\le s<t$.
                \end{enumerate}
                Jadi definisi martingale terpenuhi.
        \end{enumerate}

  \item Ambil fungsi
        \[
          F(x)=\ln x,\qquad x>0.
        \]
        Maka
        \[
          F'(x)=\frac{1}{x},\qquad F''(x)=-\frac{1}{x^2}.
        \]
        Karena $Y(t)=F(S(t))=\ln S(t)$ dan
        \[
          dS(t)=\mu S(t)\,dt+\sigma S(t)\,dW(t),
        \]
        maka menurut rumus It\^o,
        \[
          dY(t)=F'(S(t))\,dS(t)+\frac{1}{2}F''(S(t))(dS(t))^2.
        \]
        Selanjutnya,
        \[
          (dS(t))^2=(\sigma S(t)\,dW(t))^2=\sigma^2 S(t)^2\,dt,
        \]
        karena $(dW(t))^2=dt$ dan suku orde lebih tinggi diabaikan. Jadi
        \begin{align*}
          dY(t)
           & =\frac{1}{S(t)}\bigl(\mu S(t)\,dt+\sigma S(t)\,dW(t)\bigr)
          +\frac{1}{2}\left(-\frac{1}{S(t)^2}\right)\sigma^2 S(t)^2\,dt\\
           & =\mu\,dt+\sigma\,dW(t)-\frac{1}{2}\sigma^2\,dt\\
           & =\left(\mu-\frac{\sigma^2}{2}\right)dt+\sigma\,dW(t).
        \end{align*}
        Integralkan dari $0$ sampai $t$:
        \begin{align*}
          Y(t)-Y(0)
           & =\int_0^t \left(\mu-\frac{\sigma^2}{2}\right)\,ds
          +\sigma\int_0^t dW_s\\
           & =\left(\mu-\frac{\sigma^2}{2}\right)t+\sigma W_t.
        \end{align*}
        Dengan demikian,
        \[
          Y(t)=Y(0)+\left(\mu-\frac{\sigma^2}{2}\right)t+\sigma W_t.
        \]

  \item Karena $|f'(x)|\le M$ dan $|g'(x)|\le M$ untuk semua $x\in\R$, maka dengan teorema
        nilai tengah diperoleh, untuk setiap $x,y\in\R$,
        \[
          |f(x)-f(y)|\le M|x-y|,
          \qquad
          |g(x)-g(y)|\le M|x-y|.
        \]
        Jadi $f$ dan $g$ memenuhi kondisi \textit{global Lipschitz}.

        Selanjutnya, masih dari teorema nilai tengah,
        \begin{align*}
          |f(x)|
           & \le |f(x)-f(0)|+|f(0)|
          \le M|x|+|f(0)|,\\
          |g(x)|
           & \le |g(x)-g(0)|+|g(0)|
          \le M|x|+|g(0)|.
        \end{align*}
        Jadi terdapat konstanta $C>0$ sehingga
        \[
          |f(x)|+|g(x)|\le C(1+|x|),\qquad x\in\R,
        \]
        yaitu $f$ dan $g$ juga memenuhi kondisi \textit{linear growth}.

        Karena $\xi_0$ adalah variabel acak $\F_0$-terukur dan terintegral kuadrat, maka semua
        hipotesis teorema eksistensi dan ketunggalan untuk persamaan diferensial stokastik
        \[
          d\xi(t)=f(\xi(t))\,dt+g(\xi(t))\,dW(t),\qquad \xi(0)=\xi_0
        \]
        terpenuhi.

        Oleh karena itu, masalah nilai awal tersebut mempunyai solusi yang ada untuk semua
        $t\ge 0$ dan solusi itu tunggal (unik hampir pasti). Dengan kata lain, terdapat tepat
        satu proses It\^o $\xi(t)$ yang memenuhi
        \[
          \xi(t)=\xi_0+\int_0^t f(\xi(s))\,ds+\int_0^t g(\xi(s))\,dW_s,
          \qquad t\ge 0.
        \]
        Jadi terbukti bahwa masalah nilai awal tersebut mempunyai solusi tunggal.
\end{enumerate}

\end{document}
